\documentclass[aspectratio=169,10pt,english]{beamer}

\input{../../../_preambulo_beamer.tex}
\input{../../../_comandos_beamer.tex}

\selectlanguage{english}

\title{MA1001B - Statistical Modeling for Decision Making}
\subtitle{Lesson 04: Anomaly Detection and Boxplots}
\author{}
\institute{Tecnol\'ogico de Monterrey}
\date{}

\begin{document}

\begin{frame}[plain]
  \vspace{1.2cm}
  {\Large\bfseries MA1001B - Statistical Modeling for Decision Making\par}
  \vspace{0.7cm}
  {\large Lesson 04: Anomaly Detection and Boxplots\par}
  \vspace{0.7cm}
  {\color{TecRojo}\rule{\textwidth}{1.2pt}}
  \vspace{0.8cm}

  MA1001B Faculty\\[0.3cm]
  Term\\[0.3cm]
  Tecnol\'ogico de Monterrey
\end{frame}

\begin{frame}{The Screening Question}
  \begin{alertblock}{Guiding question}
    Which processing times deserve investigation, and what can a statistical
    flag actually tell us?
  \end{alertblock}

  \vspace{0.5em}
  By the end of this lesson, you can:
  \begin{itemize}
    \item state when the empirical rule is appropriate;
    \item apply Chebyshev's inequality without a shape assumption;
    \item calculate and interpret 1.5-IQR boxplot fences;
    \item separate a screening flag from a conclusion about its cause.
  \end{itemize}

  \vfill
  \scriptsize All 600 order records are synthetic. No real company data are used.
\end{frame}

\begin{frame}{Three Distance Rules}
  \small
  \begin{center}
    \begin{tabular}{p{0.22\textwidth}p{0.33\textwidth}p{0.34\textwidth}}
      \toprule
      \textbf{Rule} & \textbf{Condition} & \textbf{What it provides} \\
      \midrule
      Empirical rule & Bell-shaped and symmetric data & Approximate coverage within $1s$, $2s$, and $3s$ \\
      Chebyshev & Any distribution, with $k>1$ & Lower bound $1-1/k^2$ within $ks$ \\
      Boxplot fences & Ordered quantitative data & Values beyond the 1.5-IQR fences \\
      \bottomrule
    \end{tabular}
  \end{center}

  \vspace{0.7em}
  \begin{exampleblock}{Shared limit}
    Each rule describes statistical position. Record quality, operational
    context, and business importance require additional evidence.
  \end{exampleblock}
\end{frame}

\begin{frame}{Step 1: Empirical Rule Under Its Condition}
  \begin{columns}[T]
    \begin{column}{0.37\textwidth}
      \small
      Bell-shaped synthetic times:

      \vspace{0.4em}
      $n=600$ orders\\
      Mean $=44.87$ min\\
      Sample $s=4.87$ min

      \vspace{0.7em}
      \begin{tabular}{lr}
        \toprule
        \textbf{Interval} & \textbf{Observed} \\
        \midrule
        Within $1s$ & 69.17\% \\
        Within $2s$ & 94.83\% \\
        Within $3s$ & 99.50\% \\
        \bottomrule
      \end{tabular}

      \vspace{0.5em}
      These values support the approximate 68\%, 95\%, and more than 99\%
      pattern for this sample.
    \end{column}
    \begin{column}{0.63\textwidth}
      \centering
      \includegraphics[width=\textwidth]{../figures/anomaly_detection_01_empirical_rule.png}
    \end{column}
  \end{columns}
\end{frame}

\begin{frame}{Step 2: A Shape-Free Lower Bound}
  \begin{columns}[T]
    \begin{column}{0.39\textwidth}
      \small
      Queue delays create a right tail:

      \vspace{0.3em}
      59 delayed orders\\
      Mean $=45.79$ min\\
      Sample $s=5.84$ min\\
      Skewness $=0.74$

      \vspace{0.5em}
      \[
        \text{Minimum within }ks=1-\frac{1}{k^2}
      \]

      \begin{tabular}{lrr}
        \toprule
        & \textbf{Bound} & \textbf{Observed} \\
        \midrule
        Within $2s$ & 75.00\% & 95.33\% \\
        Within $3s$ & 88.89\% & 98.67\% \\
        \bottomrule
      \end{tabular}
    \end{column}
    \begin{column}{0.61\textwidth}
      \centering
      \includegraphics[width=\textwidth]{../figures/anomaly_detection_02_chebyshev.png}
    \end{column}
  \end{columns}
\end{frame}

\begin{frame}{Step 3: Boxplot Investigation Candidates}
  \small
  \begin{columns}[T]
    \begin{column}{0.36\textwidth}
      \[
        \mathrm{IQR}=Q_3-Q_1=6.71
      \]
      $Q_1=42.19$ min\\
      $Q_3=48.89$ min\\
      Lower fence $=32.13$ min\\
      Upper fence $=58.96$ min

      \vspace{0.8em}
      \begin{alertblock}{Verified screening result}
        The rule flags 24 orders. It captures all four injected extremes and
        also flags 20 other observations.
      \end{alertblock}
    \end{column}
    \begin{column}{0.64\textwidth}
      \centering
      \includegraphics[width=\textwidth]{../figures/anomaly_detection_03_boxplot_flags.png}
    \end{column}
  \end{columns}
\end{frame}

\begin{frame}{From Flag to Decision}
  \small
  \begin{center}
    \begin{tabular}{p{0.21\textwidth}p{0.69\textwidth}}
      \toprule
      \textbf{Stage} & \textbf{Evidence needed} \\
      \midrule
      Statistical screen & A stated rule, reproducible calculation, and preserved original value \\
      Record check & Correct unit, timestamp, identifier, and transformation history \\
      Process context & Related events, operating conditions, and comparable records \\
      Resolution & Documented reason to retain, correct, exclude, or model separately \\
      \bottomrule
    \end{tabular}
  \end{center}

  \vspace{0.8em}
  \begin{exampleblock}{What the 20 other flags mean}
    They satisfy the same 1.5-IQR rule and lack the known synthetic construction
    label. Their causes remain unknown. Calling them errors or false positives
    would require evidence that the statistical screen does not contain.
  \end{exampleblock}
\end{frame}

\begin{frame}{Concept Check}
  \begin{enumerate}
    \item Why would the empirical rule be questionable for the right-tailed
      sample in Step 2?
    \item Calculate the Chebyshev lower bound for $k=2$ and interpret it.
    \item Using $Q_1=42.19$ and $Q_3=48.89$, reproduce the upper fence.
    \item What evidence would you seek before changing a flagged record?
  \end{enumerate}

  \vspace{0.5em}
  \begin{block}{Connection to Lesson 05}
    Set notation and sample spaces provide a precise language for events. They
    will help us describe which records satisfy a rule and how events combine.
  \end{block}

  \vfill
  \scriptsize Source: OpenStax, \textit{Introductory Business Statistics 2e},
  Chapter 2, Sections 2.1, 2.2, and 2.7. CC BY-NC-SA 4.0.
\end{frame}

\appendix



\end{document}
